\textbf{Phần 1. Trắc nghiệm lựa chọn}
\setcounterpageref{ex}{0}

\begin{ex} %[3H5N35]
$1\text{ m} = \dots\text{ mm}$.
\choice
{\True $1000\text{ mm}$}
{$10\text{ mm}$}
{$100\text{ mm}$}
{$10000\text{ mm}$}
\loigiai{
$1\text{ m} = 1000\text{ mm}$.\\
\textbf{Đáp án của câu hỏi là A}
}
\end{ex}

\begin{ex} %[3H5N35]
$1\text{ kg} = \dots\text{ g}$.
\choice
{$100\text{ g}$}
{\True $1000\text{ g}$}
{$10\text{ g}$}
{$10000\text{ g}$}
\loigiai{
$1\text{ kg} = 1000\text{ g}$.\\
\textbf{Đáp án của câu hỏi là B}
}
\end{ex}

\begin{ex} %[3H5H35]
$1\text{ l} = \dots\text{ ml}$.
\choice
{\True $1000\text{ ml}$}
{$100\text{ ml}$}
{$10\text{ ml}$}
{$10000\text{ ml}$}
\loigiai{
$1\text{ l} = 1000\text{ ml}$.\\
\textbf{Đáp án của câu hỏi là A}
}
\end{ex}

\begin{ex} %[3H5H35]
Tính: $25\text{ mm} \times 4 = \dots$
\choice
{$80\text{ mm}$}
{$120\text{ mm}$}
{\True $100\text{ mm}$}
{$90\text{ mm}$}
\loigiai{
$25 \times 4 = 100\text{ mm}$.\\
\textbf{Đáp án của câu hỏi là C}
}
\end{ex}

\begin{ex} %[3H5H35]
Tính: $800\text{ g} : 4 = \dots$
\choice
{$150\text{ g}$}
{$250\text{ g}$}
{$300\text{ g}$}
{\True $200\text{ g}$}
\loigiai{
$800 : 4 = 200\text{ g}$.\\
\textbf{Đáp án của câu hỏi là D}
}
\end{ex}

\begin{ex} %[3H5H35]
Nhiệt độ ngoài trời hôm nay là $34^\circ\text{C}$, ngày mai dự báo giảm $5^\circ\text{C}$. Nhiệt độ ngày mai là
\choice
{$28^\circ\text{C}$}
{\True $29^\circ\text{C}$}
{$30^\circ\text{C}$}
{$31^\circ\text{C}$}
\loigiai{
$34 - 5 = 29^\circ\text{C}$.\\
\textbf{Đáp án của câu hỏi là B}
}
\end{ex}

\begin{ex} %[3H5V35]
Một hình tam giác có 3 cạnh dài lần lượt là $25\text{ mm}, 35\text{ mm}, 40\text{ mm}$. Chu vi hình tam giác đó là
\choice
{$90\text{ mm}$}
{$110\text{ mm}$}
{\True $100\text{ mm}$ (hay $10\text{ cm}$)}
{$80\text{ mm}$}
\loigiai{
$25 + 35 + 40 = 100\text{ mm}$.\\
\textbf{Đáp án của câu hỏi là C}
}
\end{ex}

\begin{ex} %[3H5V35]
Có 3 chai nước ngọt, mỗi chai $330\text{ ml}$. Cả 3 chai có dung tích là
\choice
{$960\text{ ml}$}
{$1000\text{ ml}$}
{$900\text{ ml}$}
{\True $990\text{ ml}$}
\loigiai{
$330 \times 3 = 990\text{ ml}$.\\
\textbf{Đáp án của câu hỏi là D}
}
\end{ex}

\begin{ex} %[3H5V35]
Một con gà nặng $1\text{ kg } 800\text{ g}$, một con vịt nặng $1\text{ kg } 500\text{ g}$. Cả hai con nặng tất cả là
\choice
{\True $3300\text{ g}$}
{$3100\text{ g}$}
{$3500\text{ g}$}
{$3200\text{ g}$}
\loigiai{
$1800 + 1500 = 3300\text{ g}$.\\
\textbf{Đáp án của câu hỏi là A}
}
\end{ex}

\begin{ex} %[3H5V35]
Đoạn thẳng $AB$ dài $14\text{ cm}$. Điểm $M$ là trung điểm của $AB$. Đoạn thẳng $AM$ dài bao nhiêu mi-li-mét?
\choice
{$60\text{ mm}$}
{$80\text{ mm}$}
{\True $70\text{ mm}$}
{$50\text{ mm}$}
\loigiai{
$AM = 14 : 2 = 7\text{ cm} = 70\text{ mm}$.\\
\textbf{Đáp án của câu hỏi là C}
}
\end{ex}

\begin{ex} %[3H5V35]
Một hộp trà nặng $450\text{ g}$, sau khi dùng bớt $\dfrac{1}{3}$ lượng trà thì lượng trà đã dùng là
\choice
{$120\text{ g}$}
{$180\text{ g}$}
{$200\text{ g}$}
{\True $150\text{ g}$}
\loigiai{
$450 : 3 = 150\text{ g}$.\\
\textbf{Đáp án của câu hỏi là D}
}
\end{ex}

\begin{ex} %[3H5C35]
Một can đựng đầy $5\text{ l}$ nước. Người ta đã lấy ra 4 chai nước, mỗi chai $800\text{ ml}$. Lượng nước còn lại trong can là
\choice
{$1600\text{ ml}$}
{\True $1800\text{ ml}$}
{$2000\text{ ml}$}
{$1500\text{ ml}$}
\loigiai{
$5\text{ l} = 5000\text{ ml}$.\\
Đã lấy: $800 \times 4 = 3200\text{ ml}$.\\
Còn lại: $5000 - 3200 = 1800\text{ ml}$.\\
\textbf{Đáp án của câu hỏi là B}
}
\end{ex}

\textbf{Phần 2. Trắc nghiệm đúng sai}
\setcounterpageref{ex}{0}

\begin{ex} %[3H5H35]
Xét tính đúng sai của các phép đổi và tính toán:
\choiceTF
{\True $2\text{ cm } 5\text{ mm} = 25\text{ mm}$.}
{\True $1\text{ l } 250\text{ ml} = 1250\text{ ml}$.}
{$1\text{ kg } 50\text{ g} = 1500\text{ g}$.}
{\True $1\text{ kg } 50\text{ g} = 1050\text{ g}$.}
\loigiai{
\begin{itemize}
    \item $2\text{ cm } 5\text{ mm} = 25\text{ mm}$: Đúng.
    \item $1\text{ l } 250\text{ ml} = 1250\text{ ml}$: Đúng.
    \item $1\text{ kg } 50\text{ g} = 1050\text{ g} \neq 1500\text{ g}$. Ý 3 Sai, Ý 4 Đúng.
\end{itemize}
\begin{itemize}
    \item a) \textbf{Đúng}.
    \item b) \textbf{Đúng}.
    \item c) \textbf{Sai}.
    \item d) \textbf{Đúng}.
\end{itemize}
}
\end{ex}

\begin{ex} %[3H5V35]
Mẹ mua $1\text{ kg}$ đường. Ngày thứ nhất dùng $250\text{ g}$, ngày thứ hai dùng $350\text{ g}$. Xét tính đúng sai:
\choiceTF
{\True Cả 2 ngày mẹ dùng hết $600\text{ g}$ đường.}
{\True Lượng đường còn lại sau 2 ngày là $400\text{ g}$.}
{Ngày thứ hai dùng ít hơn ngày thứ nhất $100\text{ g}$.}
{\True Ngày thứ hai dùng nhiều hơn ngày thứ nhất $100\text{ g}$.}
\loigiai{
\begin{itemize}
    \item Đã dùng: $250 + 350 = 600\text{ g}$. Đúng.
    \item Còn lại: $1000 - 600 = 400\text{ g}$. Đúng.
    \item Ngày 2 dùng $350 > 250$ là nhiều hơn $100\text{ g}$. Ý 3 Sai, Ý 4 Đúng.
\end{itemize}
\begin{itemize}
    \item a) \textbf{Đúng}.
    \item b) \textbf{Đúng}.
    \item c) \textbf{Sai}.
    \item d) \textbf{Đúng}.
\end{itemize}
}
\end{ex}

\textbf{Phần 3. Trả lời ngắn}
\setcounterpageref{ex}{0}

\begin{ex} %[3H5H35]
Tính: $48\text{ mm} : 6$.
\shortans{8}
\loigiai{
\textbf{Đáp số:} $8$\\
$48 : 6 = 8\text{ mm}$.
}
\end{ex}

\begin{ex} %[3H5H35]
Tính: $250\text{ g} + 450\text{ g}$.
\shortans{700}
\loigiai{
\textbf{Đáp số:} $700$\\
$250 + 450 = 700\text{ g}$.
}
\end{ex}

\begin{ex} %[3H5V35]
Tính: $1000\text{ ml} - 650\text{ ml}$.
\shortans{350}
\loigiai{
\textbf{Đáp số:} $350$\\
$1000 - 650 = 350\text{ ml}$.
}
\end{ex}

\begin{ex} %[3H5V35]
Nhiệt độ phòng lúc sáng là $24^\circ\text{C}$, lúc trưa là $31^\circ\text{C}$. Hỏi nhiệt độ đã tăng bao nhiêu độ C?
\shortans{7}
\loigiai{
\textbf{Đáp số:} $7$\\
$31 - 24 = 7^\circ\text{C}$.
}
\end{ex}

\begin{ex} %[3H5V35]
Tính chu vi hình vuông có cạnh dài $15\text{ mm}$ theo mi-li-mét.
\shortans{60}
\loigiai{
\textbf{Đáp số:} $60$\\
$15 \times 4 = 60\text{ mm}$.
}
\end{ex}

\begin{ex} %[3H5C35]
Một gói bột canh nặng $180\text{ g}$. Có 5 gói bột canh như thế chia đều cho 3 người. Mỗi người nhận được bao nhiêu gam bột canh?
\shortans{300}
\loigiai{
\textbf{Đáp số:} $300$\\
Tổng: $180 \times 5 = 900\text{ g}$.\\
Mỗi người: $900 : 3 = 300\text{ g}$.
}
\end{ex}
